%=====================================================================
\chapter{The Mathematics This Course Assumes}\label{app:maths}
%=====================================================================
\normalfigures

Everything here was taught in \emph{Discrete Structures} or \emph{Calculus}.
It is collected so that a chapter can point at it rather than re-derive it, and
each section names the chapters that need it.

\section{Logarithms}

\begin{definitionbox}[Notation]
Throughout this book:
\[ \lg n = \log_{2} n, \qquad \ln n = \log_{e} n, \qquad
   \log n = \log_{10} n. \]
When a logarithm appears inside $\bigO{\cdot}$ the base does not matter, since
$\log_{a} n = \log_{b} n / \log_{b} a$ and the divisor is a constant. That is
why \chref{master} writes $\bigTh{n\lg n}$ and never specifies a base for the
bound.
\end{definitionbox}

\begin{conceptbox}[Identities used in this book]
\rowcolors{2}{white}{lightbg}
\begin{longtable}{@{}L{6.2cm}L{8.3cm}@{}}
\toprule
\rowcolor{primary!15}
\textbf{Identity} & \textbf{Used in} \\
\midrule
\endhead
$\log(ab) = \log a + \log b$ & \chref{lowerbound}, summing $\lg(n!)$ \\
$\log(a/b) = \log a - \log b$ & \chref{recurrences}, $\lg(n/2) = \lg n - 1$ \\
$\log(a^{b}) = b \log a$ & \chref{cases}, $\lg(10^{7}) = 7\lg 10$ \\
$a^{\log_{b} n} = n^{\log_{b} a}$ & \chref{master} --- this is where the
  watershed function comes from \\
$\log_{a} n = \dfrac{\log_{b} n}{\log_{b} a}$ & change of base; why the base
  vanishes inside $\bigO{\cdot}$ \\
$\lg n \approx 3.322 \log n$ & quick mental conversion \\
\bottomrule
\end{longtable}

\smallskip
\textbf{Worth memorising:} $\lg 1000 \approx 10$, $\lg 10^{6} \approx 20$,
$\lg 10^{9} \approx 30$. A thousand is ten doublings. Nearly every
back-of-envelope estimate in this book uses one of these three.
\end{conceptbox}

\section{Summations}

\begin{conceptbox}[The six that appear]
\rowcolors{2}{white}{lightbg}
\begin{longtable}{@{}L{5.6cm}L{3.2cm}L{5.5cm}@{}}
\toprule
\rowcolor{primary!15}
\textbf{Sum} & \textbf{Closed form} & \textbf{Where} \\
\midrule
\endhead
$\displaystyle\sum_{k=1}^{n} k$ & $\dfrac{n(n+1)}{2}$ &
  insertion sort's worst case (\chref{cases}); quicksort's
  (\chref{quicksort}) \\
$\displaystyle\sum_{k=1}^{n} k^{2}$ & $\dfrac{n(n+1)(2n+1)}{6}$ &
  the cubic max-subarray (\chref{bruteforce}) \\
$\displaystyle\sum_{k=0}^{n} x^{k}$ & $\dfrac{x^{n+1}-1}{x-1}$ &
  geometric; case 3 of the master theorem \\
$\displaystyle\sum_{k=0}^{\infty} x^{k}$ ($|x|<1$) & $\dfrac{1}{1-x}$ &
  why table doubling is linear (\chref{amortised}) \\
$\displaystyle\sum_{k=0}^{\infty} \dfrac{k}{2^{k}}$ & $2$ &
  why \algname{Build-Heap} is $\bigTh{n}$ (\chref{structures}) \\
$H_{n} = \displaystyle\sum_{k=1}^{n} \dfrac{1}{k}$ & $\ln n + \gamma + \bigO{1/n}$ &
  quicksort's expected comparisons (\chref{quicksort}) \\
\bottomrule
\end{longtable}

\smallskip
$\gamma \approx 0.5772$ is the Euler--Mascheroni constant. The only fact needed
about $H_{n}$ in this book is $H_{n} = \bigTh{\lg n}$.
\end{conceptbox}

\begin{worked}[The Geometric Series, Where It Earns Its Place]
\chref{amortised} needs the total copying cost of $n$ appends with doubling:
\[ 1 + 2 + 4 + \cdots + 2^{\lfloor \lg n \rfloor}
   \;=\; \sum_{k=0}^{\lfloor \lg n \rfloor} 2^{k}
   \;=\; 2^{\lfloor \lg n \rfloor + 1} - 1 \;<\; 2n. \]
\textbf{The point is that the \emph{last} term dominates.} The sum of all
earlier doublings is less than the final one, which is why amortised analysis
works for any growth factor $> 1$: the series is geometric and converges.

\smallskip
\chref{master}'s case 3 is the same observation at the other end --- there the
\emph{first} term dominates, because the ratio is less than 1.
\end{worked}

\section{Induction}

\begin{conceptbox}[The template]
To prove $P(n)$ for all $n \ge n_{0}$:
\begin{tightnum}
  \item \textbf{Base case.} Prove $P(n_{0})$.
  \item \textbf{Inductive step.} Assume $P(k)$ for some $k \ge n_{0}$ (or for
        all $n_{0} \le k < n$, for \emph{strong} induction), and prove
        $P(k+1)$.
\end{tightnum}

\smallskip
\textbf{Where it appears:} \chref{invariants}'s loop invariants are induction
on the iteration count, with a third obligation added; \chref{recurrences}'s
substitution method is induction on $n$; and \chref{structures}'s heap bounds
are induction on the height.

\smallskip
\textbf{Strong induction} is needed whenever the recurrence refers to more than
one smaller value --- $T(n) = 1 + \sum_{j<n} T(j)$ in \chref{dpone}, for
instance.
\end{conceptbox}

\begin{alertbox}[The two mistakes]
\textbf{1. Forgetting the base case.} The step alone proves nothing: ``if
$P(k)$ then $P(k+1)$'' is satisfied by statements that are false for every $n$.

\smallskip
\textbf{2. Assuming what is to be proved.} The hypothesis is $P(k)$, for a
\emph{specific} $k$; concluding $P(k+1)$ must use it rather than restate it.
\chref{recurrences}'s strengthening trick is the legitimate version of
``assume more'': you prove a \emph{stronger} statement, so the hypothesis gives
you more to work with.
\end{alertbox}

\section{Probability}

\begin{conceptbox}[What \chref{quicksort} and \chref{hashing} need]
\textbf{Expectation.} For a random variable $X$ taking value $x$ with
probability $p(x)$,
\[ \mathbb{E}[X] = \sum_{x} x\,p(x). \]

\smallskip
\textbf{Linearity of expectation}, which is the single most useful fact:
\[ \mathbb{E}[X + Y] = \mathbb{E}[X] + \mathbb{E}[Y], \]
\textbf{whether or not $X$ and $Y$ are independent.} This is why
\chref{quicksort}'s proof can sum over all $\binom{n}{2}$ pairs without
worrying that the comparisons are correlated --- and they are heavily
correlated.

\smallskip
\textbf{Indicator variables.} $X_{A} = 1$ if event $A$ occurs and 0 otherwise.
Then $\mathbb{E}[X_{A}] = \Pr[A]$, so counting an expected number of
occurrences becomes summing probabilities. \chref{cases}'s inversion count and
\chref{quicksort}'s comparison count are both this technique.
\end{conceptbox}

\begin{worked}[Indicators and Linearity, in One Example]
\chref{cases} needs the expected number of inversions in a uniform random
permutation of $n$ elements.

\smallskip
\textbf{Step 1.} For each pair $(i,j)$ with $i<j$, let $X_{ij} = 1$ if the pair
is inverted. The total is $X = \sum_{i<j} X_{ij}$.

\smallskip
\textbf{Step 2.} $\mathbb{E}[X_{ij}] = \Pr[\text{inverted}] = \tfrac12$, since
the two orderings are equally likely.

\smallskip
\textbf{Step 3.} By linearity,
\[ \mathbb{E}[X] = \sum_{i<j} \tfrac12 = \tfrac12\binom{n}{2}
   = \frac{n(n-1)}{4}. \]

\smallskip
\textbf{Step 4.} Note what was \emph{not} needed: independence. The $X_{ij}$
are strongly dependent --- if $a<b$ and $b<c$ then $a<c$ is forced --- and
linearity does not care. Measured in \chref{cases}: 0.54, 0.56, 0.54, 0.50,
0.49 of the worst case, against the predicted 0.5.
\end{worked}

\section{Stirling's Approximation}

\begin{definitionbox}[Stirling]
\[ n! = \sqrt{2\pi n}\,\Bigl(\frac{n}{e}\Bigr)^{n}
        \Bigl(1 + \Theta\bigl(\tfrac1n\bigr)\Bigr), \]
and the two consequences this book uses:
\[ n! > \Bigl(\frac{n}{e}\Bigr)^{n}, \qquad
   \lg(n!) = n\lg n - n\lg e + \bigTh{\lg n}
           = n \lg n - 1.4427\,n + \bigTh{\lg n}. \]
\end{definitionbox}

\begin{alertbox}[Use the sharp form, not $n\lg n$]
\chref{lowerbound} measured $\lg(n!)$ against both:

\rowcolors{2}{white}{lightbg}
\begin{center}
\begin{tabular}{@{}rrrr@{}}
\toprule
\rowcolor{primary!15}
$n$ & $\lg(n!)$ & $n\lg n$ & $n\lg n - 1.443n$ \\
\midrule
1{,}000     &       8{,}529 &       9{,}966 &       8{,}523 \\
1{,}000{,}000 & 18{,}488{,}885 & 19{,}931{,}569 & 18{,}488{,}874 \\
\bottomrule
\end{tabular}
\end{center}

\smallskip
The crude $n\lg n$ is 17\% high at $n = 1{,}000$ and still 8\% high at a
million. The sharp form is correct to six decimal places. Both are
$\bigTh{n\lg n}$, and only one is usable for arithmetic.
\end{alertbox}

\section{Floors, Ceilings and Useful Inequalities}

\begin{conceptbox}[Facts used without comment in the text]
\rowcolors{2}{white}{lightbg}
\begin{longtable}{@{}L{6.0cm}L{8.5cm}@{}}
\toprule
\rowcolor{primary!15}
\textbf{Fact} & \textbf{Where it is used} \\
\midrule
\endhead
$x - 1 < \lfloor x \rfloor \le x \le \lceil x \rceil < x+1$ &
  \chref{recurrences}, dropping floors \\
$\lfloor n/2 \rfloor + \lceil n/2 \rceil = n$ &
  merge sort's split (\chref{mergesort}) \\
$1 + x \le e^{x}$ for all real $x$ &
  probability bounds \\
$\binom{n}{k} \le \left(\dfrac{en}{k}\right)^{k}$ &
  counting arguments \\
$\sum_{u \in V} \deg(u) = 2|E|$ &
  every bound in Part~V (\chref{graphs}) \\
a binary tree of height $h$ has $\le 2^{h}$ leaves &
  the sorting lower bound (\chref{lowerbound}) \\
a complete binary tree on $n$ nodes has height $\lfloor \lg n \rfloor$ &
  heaps (\chref{structures}) \\
\bottomrule
\end{longtable}
\end{conceptbox}

\section{Where Each Topic Is Needed}

\rowcolors{2}{white}{lightbg}
\begin{longtable}{@{}L{4.6cm}L{9.9cm}@{}}
\toprule
\rowcolor{primary!15}
\textbf{Topic} & \textbf{Chapters that need it} \\
\midrule
\endhead
Logarithm identities & 1, 2, 3, 5, 6, 10 --- and $a^{\log_b n} = n^{\log_b a}$
  is the key step in 6 \\
Arithmetic series & 2, 7, 9, 17 \\
Geometric series & 6, 11, 17 \\
$\sum k/2^{k} = 2$ & 12 only, and it is the whole of
  \algname{Build-Heap}'s bound \\
Harmonic numbers & 9 \\
Induction & 4, 5, 10, 12, 14, 15, 20 \\
Strong induction & 15 \\
Linearity of expectation & 2, 9, 13 \\
Indicator variables & 2, 9 \\
Stirling & 10 \\
L'H\^opital's rule & 3, for $\lg n = o(n^{\varepsilon})$ \\
\bottomrule
\end{longtable}

\begin{alertbox}[If any row is unfamiliar]
Read that row's chapters' prerequisite boxes first --- each names what it
assumes. The three most load-bearing items in the whole book are
\textbf{linearity of expectation} (\chref{quicksort}'s proof collapses without
it), \textbf{$a^{\log_{b} n} = n^{\log_{b} a}$} (\chref{master} has no
watershed function without it), and \textbf{the geometric series}
(\chref{amortised} has no result without it).
\end{alertbox}
